\section{Closed-loop inspection and observation rules}
\label{sec:agent}
\subsection{Task, state, and evidence}
Each inspection episode has a question $q$, a geographic region of interest (ROI), a task type (presence, damage, count, or spatial relation), and a reference answer $y_q$. At step $t$, the simplified vehicle state $s_t=(x_t,y_t,h_t)$ determines a rendered image $I_t=\mathcal{R}(\mathcal{M},s_t)$ from scene $\mathcal{M}$. The agent also receives structured evidence $z_t$ from the fixed perception backend and registered environment state. It may search, request a task-appropriate reobservation, answer, abstain, or stop. A reobservation is counted only when a movement executes and a new view is rendered; a repeated answer on an unchanged view is not an observation action.

Figure~\ref{fig:agent-loop} summarizes the agent loop. The benchmark retains the intermediate answer, decision, movement, evidence, generation-call keys, and terminal outcome. It scores the terminal normalized answer against $y_q$, with abstentions counted as unsuccessful answers. This is an episode-level question-answering endpoint rather than an average of building classifications.

\begin{figure}[htbp]
\centering
\resizebox{\linewidth}{!}{\begin{tikzpicture}[
  font=\footnotesize,
  >=Latex,
  state/.style={draw=blue!58!black,rounded corners=2pt,fill=blue!3,align=center,
    text width=3.25cm,minimum height=.88cm,inner sep=4pt},
  decision/.style={diamond,draw=orange!80!black,fill=orange!7,aspect=2.15,
    align=center,text width=2.4cm,inner sep=1.5pt},
  terminal/.style={draw=green!45!black,rounded corners=2pt,fill=green!6,align=center,
    text width=2.7cm,minimum height=.8cm},
  arr/.style={->,thick,draw=blue!58!black},
  control/.style={draw=black!45,dashed,rounded corners=2pt,fill=gray!4,align=center,
    text width=3.2cm,inner sep=4pt}
]
\node[state] (init) {Task $q$, target ROI, initial UAV state and action budget};
\node[state,right=8mm of init] (render) {Render georeferenced observation $I_t$ and read state $s_t$};
\node[state,right=8mm of render] (infer) {Extract damage/spatial evidence $z_t$; generate Qwen candidate};
\node[decision,below=9mm of infer] (enough) {Evidence sufficient for the task?};
\node[terminal,right=9mm of enough] (answer) {Answer and report terminal outcome};
\node[decision,left=16mm of enough] (feasible) {Budget and action feasible?};
\node[terminal,left=9mm of feasible] (abstain) {Abstain or stop; log terminal outcome};
\node[state,below=10mm of feasible] (choose) {Select search, center, descend, or a composed movement};
\node[state,right=11mm of choose] (execute) {Execute action, decrement budget, and append episode log};
\node[control,below=9mm of answer] (fixed) {Fixed-view control: repeat answer generation while holding $s_t$, $I_t$, and $z_t$ constant};

\draw[arr] (init)--(render);
\draw[arr] (render)--(infer);
\draw[arr] (infer)--(enough);
\draw[arr] (enough)--node[above]{yes}(answer);
\draw[arr] (enough)--node[above]{no}(feasible);
\draw[arr] (feasible)--node[above]{no}(abstain);
\draw[arr] (feasible)--node[right]{yes}(choose);
\draw[arr] (choose)--(execute);
\draw[arr] (execute.east) -- ++(.55,0) |- (render.south);
\draw[->,thick,dashed,draw=black!55] (infer.east) -- ++(.55,0) |- (fixed.north);
\end{tikzpicture}
}
\caption{Closed-loop inspection state machine. Feasibility and budget checks distinguish a requested reobservation from an executed view change. The unchanged-view branch is available for diagnosis but is not counted as reobservation.}
\label{fig:agent-loop}
\end{figure}

\subsection{Fixed binary perception and cross-view memory}
The main experiments use the same frozen ChangeOS ResNet-34 checkpoint for every policy. It detects buildings in registered pre-/post-disaster imagery and supplies a binary no-damage/damaged probability vector. ChangeOS is an external perception tool, not a newly trained contribution of this study. Its checkpoint has prior xBD exposure; neither the perception evaluation nor the online results demonstrate transfer of the detector to unseen disaster events.

For a binary vector $p_i=(p_{i0},p_{i1})$, the damage-uncertainty score is
\begin{equation}
 H_2(p_i)=-\frac{1}{\log 2}\sum_{j=0}^{1}p_{ij}\log p_{ij}.
 \label{eq:entropy}
\end{equation}
The raw-entropy policy uses the uncalibrated binary probabilities, whereas calibrated entropy and a prediction-set trigger are supplemental comparisons. Historical four-class temperature and APS thresholds are not transferred to this experiment. For count tasks, the task-conditioned gate estimates uncertainty over the answer buckets $0,1,2,3+$ from the binary building probabilities; for spatial tasks it uses registered target geometry and a GSD-dependent bearing uncertainty. Thus a single uncertain building label is not treated as a sufficient reason to crop away a complete count or spatial ROI.

The online evidence memory matches predicted objects by geographic proximity and combines observations across views, preserving the history of observation IDs and ROI coverage. This avoids treating each new crop as an unrelated scene or automatically discarding earlier wide-view evidence. Answer generation uses a locally loaded Qwen2.5-VL-7B-Instruct model with the task image and selected structured state; the evaluated configurations use the same hybrid answer path. The model's output is normalized to the task's answer vocabulary before scoring. A deterministic rule answerer and pure-VLM paths exist as development diagnostics, but are not pooled into the seven-policy final comparison.

\subsection{Task-conditioned gate and comparators}
The task-conditioned policy (T1\_TASK) estimates whether an executable next view is worth its cost. In its implemented rule, the candidate utility is
\begin{equation}
 U_q= u_q(1-h_{t+1}/h_t)
 -0.05\,\widehat{T}_{\mathrm{move}}/60
 -(1-\widehat{C}_{\mathrm{ROI}}),
 \label{eq:taskutility}
\end{equation}
where $u_q$ is task-specific uncertainty, $\widehat{T}_{\mathrm{move}}$ is a configured-speed motion-time estimate in seconds, and $\widehat{C}_{\mathrm{ROI}}$ is predicted ROI coverage after the candidate action. This is a fixed heuristic, not a learned or calibrated estimate of expected task accuracy. The frozen gate requires $u_q\geq0.5$ and $U_q\geq0.05$; count and spatial tasks additionally require predicted coverage of at least 0.98. A damage question may center on a visible marked target and descend. Count and spatial questions retain the geographic center and descend only if the task ROI remains covered. Presence positives are answered from the wide view, while missing evidence uses the search path. All thresholds were chosen before the final run using development evidence.

The four primary comparators are A0\_HOLD (search allowed, no reobservation), A1\_RANDOM (independent seeded random trigger), A2\_ALWAYS (request whenever feasible), and A3U\_RAW\_ENTROPY (binary raw-entropy trigger). A3\_ENTROPY and A4\_CONFORMAL are supplemental. Each active policy has the same exogenous cap of two reobservations per episode, and all configurations permit up to six search steps. A0\_HOLD has a zero reobservation cap. The cap is shared, but realized action counts differ because policy requests and feasibility differ; the seven outcomes are operating points, not a randomized budget sweep.
